% 从旧数学附录迁出的应用说明；不新增算法。
\section{回文半径与在线回文计数}\label{knowledge-palindrome}
Manacher 的 odd[i] 包含中心，最长奇回文为 $[i-odd[i]+1,i+odd[i]-1]$；
even[i] 以 i 前间隙为中心，最长偶回文为 $[i-even[i],i+even[i]-1]$。
odd/even 之和是回文出现总次数（long long），不是不同回文串数。
按字节比较，空串合法，int 索引要求 $n<2^{30}$。

PAM 的 suffixes 等于 fail 节点的 suffixes 加一；hits 只计最长后缀命中，完整次数用 occurrences()。
total 最多 $n(n+1)/2$，distinct 才是不同非空回文串数。PAM 固定小写字母。
\section{整数后缀数组的边界与 SAM 次数}\label{knowledge-suffix-bounds}
SuffixArray 整数入口要求 $0\le a[i]<alphabet$，大整数先保序离散化。
sa/rk/lcp 从零编号，lcp[k] 比较 sa[k-1] 与 sa[k]，首项为零。
LCP 需遇分隔符停止时，令 remain[i] 为 i 到下一分隔符前的普通字符数：
\[
 height[k]=\min\{lcp[k],remain[sa[k-1]],remain[sa[k]]\}.
\]
两串拼接求最长公共子串只需检查异来源的相邻后缀。

SAM 状态 u 代表长度 $(len[link[u]],len[u]]$ 内同一 endpos 的子串。
occ 种子真实状态为一、克隆为零，按 len 倒序传播（不能按创建编号）。
根计数为 $n$，非空子串只看非根状态；次数乘长度用 long long。
\section{广义 SAM 点数与后缀语言}\label{knowledge-general-sam}
GeneralSAM 的状态按（文档编号，结束位置）的 endpos 集合等价。
先 add 全部小写串，再 build；根为零，distinct 为不同非空子串数，a.size() 含根。
构造后不再 add。多串 endpos 自动机不是后缀并集的最小 DFA，点数按 endpos 口径输出 a.size()。
\section{排列约束与最大流单位边分类}\label{knowledge-flow-units}
已知元素切开未知段：段内候选来自原排列对应区间，需求为未知位置数；已知元素单独成需求一的组，
首尾设虚拟锚点，保留零需求组。源连左组、右组连汇，容量为需求，每个元素是左右组间一条单位边；
流量达到 $k$ 等价于存在合法集合。
满流后在残量网络求 SCC，判断能否沿环改变某条单位边，得到必选/可选/不可选。
本题增广求流 $O(n^2)$，SCC 为 $O(n)$。
\section{点分治、启发式合并与虚树用法}\label{knowledge-tree-uses}
点分治 count\_exact 计无序不同顶点对，预处理 $O(n\log^2n)$，每次查询 $O(n\log n)$；边修改后须重建。

启发式合并：轻子树先算后清空，重子树保留；并列最高频颜色编号之和须用 long long。

虚树：显式加入根后构造，只清零本次 vertices 的 dp；逆序遍历虚边，
孩子为关键点取路径最小边权，否则取它与孩子 dp 的较小值（用 minimum，不用 distance）。
\section{模2压位消元与完整解空间}\label{knowledge-gauss-xor}
\texttt{GaussXor::solve(a,n)} 输入 $m$ 行 $n+1$ 列 01 串，末列为右端项。
自由变量全设零得特解；每次一个自由变量设一得零空间基。
全部解为特解异或任意基子集，维数 $n-rank$。
每行 $\lfloor n/64\rfloor+1$ 个字；无解与零变量有解都可能返回空特解，须查 \texttt{consistent}。
输出完整基最坏 $n^2$ 个字符；特解不保证一的个数最少。
\section{严格次小生成树与换边}\label{knowledge-second-mst}
\texttt{SecondMST} 先 Kruskal 求 MST，倍增维护树路径上两个不同的最大边权及边号。
每条非树非自环边换入：非严格删路径最大边；严格且换入边等于最大值时删第二大不同权边，无则跳过。
\texttt{tree} 删 \texttt{out}、加 \texttt{in} 得 \texttt{next}；无次小树时 \texttt{next} 为空，
断图时 \texttt{connected=false}。总权用 int128，边号 $-1$ 表示缺项。
\section{模2消元后的最少开关用法}\label{knowledge-switch-min}
每个方格一个变量，系数为自身及四邻，右端项为初始颜色。
解集为特解异或零空间基任意子集，须枚举取一最少者；无解输出 \texttt{inf}。
自由维数 $d$ 时代价 $O(n^2 2^d)$；对此四邻棋盘 $d\le n$，不可推广到任意方程组。


\Needspace{360pt}
\section{闭合集、树形图与全局割的完整用法}\label{knowledge-closure-cuts}
最大权闭合子图：收益为正权点、成本为负权点，依赖边由前者指向后者；
maximum\_closure 返回所选点号，共享成本只计一次，总额用 long long。

Arborescence::solve 点号与根 0-based，返回 optional 区分无解与负费用；只返回费用，不恢复边。
StoerWagner::solve 要求至少两点、非负对称权，side 为一侧 0-based 原顶点。
核心费用为 128 位类型。
